\chapter{Grassmannians and the space of lines in \texorpdfstring{$\mathbb{P}^3$}{P3}}
\label{grassmannians}

In this chapter we will introduce Grassmann varieties, or Grassmannians, and focus on $\mathbb{G}(1, 3)$, which is the space of lines in $\mathbb{P}^3$. We will also introduce the Plücker embedding, which is a way to embed the Grassmannian into a projective space, giving it the structure of a projective variety. For a more detailed introduction to Grassmannians, we refer to \cite{EH16}.

\section{Definition and basic properties}
A Grassmannian is a projective variety whose closed points correspond to linear subspaces of a certain dimension in an $n$-dimensional vector space. As a set, $\mathrm{Gr}(k, V)$ is the set of all $k$-dimensional linear subspaces of the vector space $V$. In \cref{sec:Plucker}, we will give it the structure of a projective variety. In our case, $V$ will be a vector space over $\mathbb{C}$. \todo[inline]{What happens when we use another field than $\mathbb{C}$? EO works over $\mathbb{Z}$ to get all fields.}
\[\mathrm{Gr}(k,n)(\mathbb{C}) = \{[L] \mid L \subseteq \mathbb{k}^n \text{ is a linear space of dimension } k\}.\]
\todo{\cite{EO25}}

A $k$-dimensional linear subspace of an $n$-dimensional vector space $V$ can also be thought of as a $(k - 1)$-dimensional linear subspace of $\mathbb{P}V \cong \mathbb{P}^{n-1}$. Due to this, we will often write $\mathrm{Gr}(k, V)$ as $\mathbb{G}(k-1, \mathbb{P}V)$. When there is no need to specify the vector space $V$, only its dimension $n$, we will write $\mathrm{Gr}(k,n) = \mathbb{G}(k-1, n-1)$.

% To give it the structure of a projective variety, we give an inclusion in a projective space, called the Plücker embedding, and show that the image is the zero locus of a certain collection of homogeneous polynomials.


\cite{EO25} A dimension $n-k$ linear subspace $L \subseteq \mathbb{k}^n$ is described as the kernel of a surjective linear map $\mathbb{k}^n \to \mathbb{k}^{n-k}$, represented by an $(n-k) \times n$ matrix $a$ of rank $n-k$:
\[L = \text{Ker}(A) \subseteq \mathbb{k}^n,\]

where 
\[A = \begin{pmatrix}
a_{1,1} & a_{1,2} & \cdots & a_{1,n} \\
a_{2,1} & a_{2,2} & \cdots & a_{2,n} \\
\vdots & \vdots & \ddots & \vdots \\
a_{(n-k),1} & a_{(n-k),2} & \cdots & a_{(n-k),n}
\end{pmatrix}.\]

Note that the matrix $A$ is not unique for a given linear subspace $L$, since multiplying $A$ on the left by an invertible $(n-k) \times (n-k)$ matrix $G$ gives another matrix $GA$ with the same kernel as $A$. Thus, we consider the equivalence relation on the set of $(n-k) \times n$ matrices of rank $n-k$ given by $A \sim A'$ if there exists an invertible $(n-k) \times (n-k)$ matrix $G$ such that $A' = GA$. The Grassmannian $\mathrm{Gr}(k,n)$ can then be identified with the set of equivalence classes of $(n-k) \times n$ matrices of rank $n-k$ under this equivalence relation.

\section{The Plücker embedding} \label{sec:Plucker}

\begin{itemize}
    \item General idea of Plucker embedding, referencing EO and EH.
    \item Example: $\mathbb{G}(1,3)$.
\end{itemize}

\section{Covering by affine spaces}

\begin{itemize}
    \item Idea of covering by affine spaces, referencing EO and EH.
    \item Example: $\mathbb{G}(1,3)$.
\end{itemize}

\begin{example}
For $n=4, k=2$, there are 6 affine spaces $\mathbb{A}^4$. They are:

\begin{align*}
    U_{1,2} &=\begin{pmatrix}1 & 0 & a_{1,3} & a_{1,4} \\0 & 1 & a_{2,3} & a_{2,4}\end{pmatrix} \\
    U_{1,3} &=\begin{pmatrix}1 & a_{1,2} & 0 & a_{1,4} \\0 & a_{2,2} & 1 & a_{2,4}\end{pmatrix} \\
    U_{1,4} &=\begin{pmatrix}1 & a_{1,2} & a_{1,3} & 0 \\0 & a_{2,2} & a_{2,3} & 1\end{pmatrix} \\
    U_{2,3} &=\begin{pmatrix}a_{1,1} & 1 & 0 & a_{1,4} \\a_{2,1} & 0 & 1 & a_{2,4}\end{pmatrix} \\
    U_{2,4} &=\begin{pmatrix}a_{1,1} & 1 & a_{1,3} & 0 \\a_{2,1} & 0 & a_{2,3} & 1\end{pmatrix} \\
\end{align*}

\end{example}
\section{Short on Schubert cycles (?)}
\todo[inline]{Decide if we should include or not. Probably not.}

